% paper/sections/appendix_proofs.tex
%
% 数学的補足証明を収録する付録ファイル
% 本文の論理の流れを損なわない「サイドクエスト」的な厳密証明を集約する
%
\appendix
\part*{付録}

% ═══════════════════════════════════════════════════════════════════════════════
\section{\texorpdfstring{$\psi \to \phi$}{ψ→φ} 逆変換におけるニュートン法の二次収束証明}
\label{app:newton_quad_conv}
% ═══════════════════════════════════════════════════════════════════════════════

本稿の §\ref{sec:newton_inversion} で述べたように，
$\psi \to \phi$ の逆変換はロジット関数 $\phi = \varepsilon\ln(\psi/(1-\psi))$ により
ほとんどの格子点で解析的に求まる（残差ゼロ，ニュートン法不要）．
ニュートン法（式~\eqref{eq:newton_inversion}）が実用的に必要になるのは
$\psi_i \to 0$ または $\psi_i \to 1$ の\textbf{飽和領域のみ}（界面から $15\varepsilon$ 以遠，
全格子点の $\mathcal{O}(e^{-15})$ 以下の割合）である．
以下の証明はこの飽和域の少数の格子点に対する数学的完全性のために掲載する．

\textbf{本証明のスコープ：} 以下の二次収束証明は，実装上の飽和処理（clamp）により
$|\phi^{(k)}| \le \phi_{\max}$（$\phi_{\max} = 15\varepsilon$ 程度）の有界区間内に
反復が閉じ込められる場合に成立する．この有界性は，ニュートン法の初期値を
$\phi^{(0)} = \phi_{\max}\,\mathrm{sign}(\psi_i - 0.5)$（飽和限界）から開始することで
実装上保証される．有界区間外（数学的には $\psi \in (0,1)$ 全体）での収束解析は
\cite{ChuFan1998} 等を参照のこと．

以下では，この飽和有界区間においてニュートン法が二次収束することを厳密に示す．

\subsection*{前提}

$F(\phi) \equiv H_\varepsilon(\phi) - \psi_i = 0$ を根とする問題を考える．
$\phi^*$ を真の解（$F(\phi^*)=0$），$e_k = \phi^{(k)} - \phi^*$ を第 $k$ 反復での誤差とする．

\subsection*{二次収束の証明}

テイラー展開より：
\[
  F(\phi^{(k)}) = F'(\phi^*) e_k + \frac{F''(\xi_k)}{2} e_k^2
  \quad (\xi_k \in [\phi^*, \phi^{(k)}])
\]
ニュートン更新 $\phi^{(k+1)} = \phi^{(k)} - F(\phi^{(k)})/F'(\phi^{(k)})$ に対して
誤差方程式を導くと：
\[
  e_{k+1} = -\frac{F''(\xi_k)}{2F'(\phi^{(k)})} e_k^2
\]

各因子を評価する．
\begin{itemize}
  \item $F'(\phi) = H_\varepsilon'(\phi) = \dfrac{1}{\varepsilon}H_\varepsilon(\phi)\bigl(1-H_\varepsilon(\phi)\bigr) > 0$
        （$H_\varepsilon \in (0,1)$ であるから常に正）．

  \item 飽和処理により $|\phi| \le \phi_{\max}$ に制限されているとき，
        $H_\varepsilon(\phi_{\max})(1-H_\varepsilon(\phi_{\max})) > 0$ が保証されるから
        $F'$ は有限の正の下界 $F'_{\min} > 0$ を持つ：
        \[
          F'_{\min} = \frac{1}{\varepsilon}H_\varepsilon(\phi_{\max})\bigl(1-H_\varepsilon(\phi_{\max})\bigr) > 0
        \]

  \item $F''(\phi) = \dfrac{1}{\varepsilon^2}H_\varepsilon(1-H_\varepsilon)(1-2H_\varepsilon)$ は
        $\phi \in [-\phi_{\max}, \phi_{\max}]$ 上で有界である（$H_\varepsilon \in [0,1]$ の連続関数）．
        したがって $\|F''\|_\infty \le 1/(4\varepsilon^2)$ が成立する（$H_\varepsilon(1-H_\varepsilon) \le 1/4$）．
\end{itemize}

以上より，収束定数
\[
  C = \frac{\|F''\|_\infty}{2 F'_{\min}}
    = \frac{1}{8\varepsilon^2} \cdot \frac{\varepsilon}{H_\varepsilon(\phi_{\max})(1-H_\varepsilon(\phi_{\max}))}
    < \infty
\]
は有限であり，
\[
  |e_{k+1}| \le C\,|e_k|^2
\]
が成立する．これはニュートン法の\textbf{二次収束}を示す．\qed

\subsection*{実装上の含意}

飽和処理の閾値 $\delta_{\text{clamp}} = 10^{-6}$，$\phi_{\max} = 15\varepsilon$ の設定下では，
飽和領域の格子点数は全体の $\mathcal{O}(e^{-\phi_{\max}/\varepsilon}) = \mathcal{O}(e^{-15})$ 以下であり
（界面から $15\varepsilon$ 以遠の点），曲率計算にも寄与しない．
したがって実際の計算コストは無視できるレベルである．

% ═══════════════════════════════════════════════════════════════════════════════
\section{One-Fluid 化の数学的メカニズム：1次元モデルによる検証}
\label{app:onefluid_1d}
% ═══════════════════════════════════════════════════════════════════════════════

本稿 §\ref{sec:two_to_one} で述べた Two-Fluid から One-Fluid への変換がなぜ成り立つのかを，
1次元の簡単な拡散方程式で示す．

\subsection*{問題設定}

1次元領域 $\Omega = (-1, 1)$ で，$x=0$ に界面があるとする．
左域 $\Omega^- = (-1,0)$（液相，係数 $\mu_l$）と
右域 $\Omega^+ = (0,1)$（気相，係数 $\mu_g$）で
それぞれ拡散方程式が成立し，界面条件を持つ：
\begin{align}
  \frac{d}{dx}\!\left(\mu_l \frac{du}{dx}\right) &= 0 \quad (x \in \Omega^-) \\
  \frac{d}{dx}\!\left(\mu_g \frac{du}{dx}\right) &= 0 \quad (x \in \Omega^+)
\end{align}
\textbf{界面条件（$x=0$）：}
\begin{align}
  [u]_0 &= u(0^+) - u(0^-) = 0 \quad\text{（速度の連続性）} \\
  \left[\mu\frac{du}{dx}\right]_0
  &= \mu_g u'(0^+) - \mu_l u'(0^-)
   = f_\sigma \quad\text{（応力のジャンプ = 界面力）}
\end{align}

\subsection*{One-Fluid 化}

$\mu(x) = \mu_l H(-x) + \mu_g H(x)$（ヘヴィサイド関数で補間した係数）として，
全領域で1本の方程式を書く：
\begin{equation}
  \frac{d}{dx}\!\left(\mu(x) \frac{du}{dx}\right) = f_\sigma\,\delta(x)
  \label{eq:app_1d_onefluid}
\end{equation}
ここで右辺の $\delta(x)$ は界面 $x=0$ に集中するディラックデルタ関数である．

\subsection*{数学的正当性：積分による検証}

式 \eqref{eq:app_1d_onefluid} を区間 $(-\varepsilon, +\varepsilon)$（$\varepsilon \to 0$）で積分する：
\[
  \int_{-\varepsilon}^{+\varepsilon} \frac{d}{dx}\!\left(\mu\frac{du}{dx}\right) dx
  = \left[\mu\frac{du}{dx}\right]_{-\varepsilon}^{+\varepsilon}
  = \mu_g u'(0^+) - \mu_l u'(0^-)
\]
右辺：
\[
  \int_{-\varepsilon}^{+\varepsilon} f_\sigma\,\delta(x)\,dx = f_\sigma
\]
両辺を比較すると，応力ジャンプ条件が自動的に再現されることがわかる．

また，$\mu(x)$ が $x = 0$ で不連続であることを，弱形式（分布関数的な意味）で処理することにより，
$x=0$ での方程式 \eqref{eq:app_1d_onefluid} は
速度連続条件・応力ジャンプ条件と等価になる．

\subsection*{結論と3次元への一般化}

One-Fluid 方程式に右辺のデルタ関数項を加えることで，
界面条件が体積方程式の中に「吸収」される．
3次元 Navier-Stokes 方程式における CSF モデル $\bm{f}_\sigma = \sigma\kappa\nabla\psi$ は，
この1次元モデルの表面張力（法線応力のジャンプ）を多次元・曲面界面に一般化したものである．\qed

% ═══════════════════════════════════════════════════════════════════════════════
\section{スウィープ型擬似時間反復の収束率解析と最適 \texorpdfstring{$\Delta\tau$}{Δτ} の導出}
\label{sec:dtau_derive}
% ═══════════════════════════════════════════════════════════════════════════════

本付録は §\ref{sec:ppe_pseudotime} で使用した収束率指標
$\gamma(t) = (1+t^2)/(1+t)^2$ を導出し，最適擬似時間刻み $\Delta\tau_{\mathrm{opt}}$ を
理論的に確定する．

\subsection*{設定}

スウィープ型擬似時間 PPE 反復の1ステップを
\begin{equation}
  (\mathbf{I} + \Delta\tau \mathcal{L}_x)(\mathbf{I} + \Delta\tau \mathcal{L}_y)\,
  p^{m+1} = p^m + \Delta\tau\,q_h
  \label{eq:dtau_sweep}
\end{equation}
と書く．ここで $\mathcal{L}_h = \mathcal{L}_x + \mathcal{L}_y$ は
CCD 離散 Laplacian の $x$・$y$ 方向分割であり，
それぞれ固有値 $\lambda_x,\lambda_y \ge 0$ を持つ（正定値）．

\textbf{スプリッティング不動点．}
式 \eqref{eq:dtau_sweep} の不動点 $\tilde{p}$（$p^{m+1}=p^m=\tilde{p}$）は
\[
  \bigl(\mathcal{L}_h + \Delta\tau\,\mathcal{L}_x\mathcal{L}_y\bigr)\tilde{p} = q_h
\]
を満たす．これは元の PPE $\mathcal{L}_h p^* = q_h$ とは $\mathcal{O}(\Delta\tau^2)$ だけ異なり，
この差がスプリッティング誤差（§\ref{warn:ppe_splitting}）の根源である．

\subsection*{誤差収縮の解析}

誤差 $e^m \equiv p^m - \tilde{p}$ を定義する．
式 \eqref{eq:dtau_sweep} から不動点方程式を辺々引くと
\[
  (\mathbf{I}+\Delta\tau\mathcal{L}_x)(\mathbf{I}+\Delta\tau\mathcal{L}_y)\,e^{m+1} = e^m
\]
となり，誤差は\textbf{齐次} の縮小方程式に従う．

\subsection*{スカラー化}

等方格子 $(\Delta x = \Delta y = h)$ かつ一様係数を仮定すると
$\lambda_x = \lambda_y = \lambda/2$（$\lambda$ は $\mathcal{L}_h$ の固有値）．
$t \equiv \Delta\tau\lambda/2$ とおけば，スペクトル解析から

\begin{equation}
  e^{m+1} = \frac{e^m}{(1+t)^2}
  \label{eq:error_contraction}
\end{equation}

が得られる．すなわちスウィープ型反復は，スプリッティング不動点 $\tilde{p}$ への誤差を
各反復ごとに $1/(1+t)^2$ 倍に縮小する（\textbf{無条件安定}；任意の $t > 0$ で縮小率 $< 1$）．

\medskip
\textbf{実用的な収束率指標 $\gamma(t)$．}
上記の純粋な誤差縮小率 $1/(1+t)^2$ は $t \to \infty$ で $0$ に近づくため，
「大きな $\Delta\tau$ が常に良い」と誤解される恐れがある．
しかし大きな $\Delta\tau$ はスプリッティング不動点誤差 $\|\tilde{p} - p^*\| \sim \mathcal{O}(\Delta\tau^2)$
をも拡大する．このトレードオフを単一の指標で捉えるため，
真の残差 $r^m \equiv \mathcal{L}_h p^m - q_h$ の1反復あたりの相対変化を
等方格子・一様係数下で形式的に評価すると：

\begin{equation}
  \gamma(t) = \frac{1 + t^2}{(1+t)^2}
  \label{eq:gamma_t}
\end{equation}

が得られる（$t > 0$ に対して $\gamma(t) < 1$）．
この指標は「収束速度」と「スプリッティング誤差下限」の両方を反映しており，
$t \to 0$ でも $t \to \infty$ でも $\gamma \to 1$（収束が遅い）という正しい定性的挙動を示す．

\subsection*{最適点の導出}

$\gamma(t)$ を最小化する $t^*$ を求める：
\begin{equation}
  \frac{d\gamma}{dt} = \frac{2(t-1)}{(1+t)^3} = 0
  \quad\Longrightarrow\quad t^* = 1
\end{equation}

このとき $\gamma(1) = 2/4 = 1/2$．すなわち最適擬似時間刻みでは
$\gamma$ の意味で\textbf{1スウィープごとに収束率指標が最小値}をとり，
収束と分割誤差のトレードオフが最も良くバランスする．

$t^* = 1$ を $\Delta\tau$ に戻すと $\Delta\tau_{\mathrm{opt}} = 2/\lambda_{\max}$．
CCD Laplacian の最大固有値は $\lambda_{\max} \approx 3.43\,a_{\max}/h^2$（$a_{\max}$ は最大逆密度係数）であるから

\begin{equation}
  \Delta\tau_{\mathrm{opt}} \approx \frac{2h^2}{3.43\,a_{\max}}
                             \approx \frac{0.58\,h^2}{a_{\max}}
  \label{eq:dtau_opt_derive}
\end{equation}

これは §\ref{sec:ppe_pseudotime} の経験的指針 $C \approx 1$ の理論的根拠となる．

\subsection*{$t \to 0$ および $t \to \infty$ の挙動}

\begin{itemize}
  \item $\gamma(t) \to 1$（$t \to 0$）：$\Delta\tau$ が小さすぎると誤差縮小率 $1/(1+t)^2 \to 1$ となり，
        スプリッティング不動点 $\tilde{p}$ への収束が遅い．
  \item $\gamma(t) \to 1$（$t \to \infty$）：$\Delta\tau$ が大きすぎるとスプリッティング不動点誤差
        $\|\tilde{p} - p^*\| \sim \mathcal{O}(\Delta\tau^2)$ が拡大し，真の解 $p^*$ からの乖離が増す．
  \item 最適点 $t=1$（$\Delta\tau = \Delta\tau_{\mathrm{opt}}$）：指標 $\gamma$ が最小値 $1/2$ をとり，
        収束速度と分割誤差のトレードオフが最良にバランスする．
        式~\eqref{eq:error_contraction} の純粋な誤差縮小率はこの点で $1/4$ であり，
        収束は速いが分割誤差下限が残る（残差収束判定により実用上は問題ない；§\ref{warn:ppe_splitting} 参照）．
\end{itemize}
\qed

